\subsection{左微分}

定义 9. 设 G 是李群，M 是 \(\mathbf{C}^r\) 类流形，\(f\) 是从 M 到 G 的 \(\mathbf{C}^r\) 类映射。\(f\) 的左（分别为右）微分是 M 上取值于 L(G) 的 1 阶微分形式，它把每个向量 \(u \in \mathrm{T}_m(\mathrm{M})\) 对应到元素 \(f(m)^{-1}\) . \((\mathrm{T}_m f)(u)\)（分别为 \((\mathrm{T}_m f)(u). f(m)^{-1}\)）。

本章只考虑左微分，记为 \(f^{-1}.df\)，右微分的相应结果留给读者翻译。

若 \(f\) 是恒等映射；\(G, f^{-1}.df\) 是左典范微分形式 \(\omega\)，属于 \(G\)。回到定义8的一般情形，

\[
(f ^ {- 1}. d f) _ {m} = \omega_ {f (m)} \circ \mathrm{T} _ {m} (f)
\]

因此 \(f^{-1}.df = f^{*}(\omega)\)。这说明 \(f^{-1}.df\) 属于 \(\mathbf{C}^{r - 1}\) 类。

例。(1) 若 G 是完备可赋范空间的加法群，且 \(T_{0}(E)\) 与 E 典范地识别，则 \(f^{-1}\) . df 就是《可微与解析流形》R，8.2.2 中定义的微分 df。

\begin{enumerate}
\def\labelenumi{(\arabic{enumi})}
\setcounter{enumi}{1}
\tightlist
\item
  假设 G 是与完备可赋范代数 A 相关的乘法群 A*。于是 f 可看作从 M 到 A 的映射，因而《可微与解析流形》R，8.2.2 意义下的微分 df 有定义，而《可微与解析流形》R，8.3.2 意义下的乘积 \(f^{-1}df\) 也有定义。显然，后一个形式与 f 的左微分相同。
\end{enumerate}

命题 58. 设 G 和 H 是两个李群，M 是 \(\mathbf{C}^r\) 类流形，\(f\) 是从 M 到 G 的 \(\mathbf{C}^r\) 类映射，\(h\) 是从 G 到 H 的态射。则

\[
(h \circ f) ^ {- 1}. d (h \circ f) = \mathrm{L} (h) \circ (f ^ {- 1}. d f) = (h ^ {- 1}. d h) \circ \mathrm{T} (f).
\]

对任意 \(x \in \mathbf{M}\) 及 \(u \in \mathrm{T}_x(\mathrm{M})\)，

\[
(h \circ f) ^ {- 1}. d (h \circ f) (u) = ((h \circ f) (x)) ^ {- 1}. T (h \circ f) (u).
\]

后一表达式一方面等于

\[
\begin{array}{l l} \mathrm{T} (h) (f (x) ^ {- 1}. \mathrm{T} (f) (u)) & (\S 2, \text { Proposition   5 }) \\ = \mathrm{T} _ {e} (h) ((f ^ {- 1}. d f) (u)) \end{array}
\]

另一方面等于

\[
\begin{array}{c} h (f (x)) ^ {- 1} \mathrm{T} (h) (\mathrm{T} (f) u) \\ = (h ^ {- 1}. d h) (\mathrm{T} (f) u). \end{array}
\]

命题 59. 设 G 是李群，M 是 \(\mathbf{C}^r\) 类流形，\(f\) 和 \(g\) 是从 M 到 G 的 \(\mathbf{C}^r\) 类映射，\(p\) 是从 TM 到 M 的典范满射。

\[
(i) (f g) ^ {- 1}. d (f g) = (\mathrm{Ad} \circ g \circ p) ^ {- 1} \circ (f ^ {- 1}. d f) + g ^ {- 1}. d g.
\]

\begin{enumerate}
\def\labelenumi{(\roman{enumi})}
\setcounter{enumi}{1}
\tightlist
\item
  定义 \(h(m) = f(m)^{-1}\)，其中 \(m \in M\) 任意。
\end{enumerate}

\[
h ^ {- 1}. d h = - (\mathrm{Ad} \circ f \circ p) \circ (f ^ {- 1}. d f).
\]

断言 (i) 由 § 2 第2号命题7推出。令 \(g = h\)，由 (i) 得到断言 (ii)。

推论 1. 设 \(s \in G\)，且 \(sg\) 是映射 \(x \mapsto sg(x)\)，其定义域为 \(M\)，陪域为 \(G\)。则 \((sg)^{-1}.d(sg) = g^{-1}.dg\)。

这由命题59 (i) 推出，其中取 \(f\) 为常值映射 \(x \mapsto s\)，其定义域为 \(M\)，陪域为 \(G\)。

推论 2. 若映射 \(f\) 和 \(g\) 的定义域为 \(M\)、陪域为 \(G\)，且它们具有相同的左微分，则 \(fg^{-1}\) 的切映射处处为零。若进一步 \(K\) 的特征为 0，则 \(fg^{-1}\) 局部为常值。

根据命题59，

\[
(f g ^ {- 1}) ^ {- 1}. d (f g ^ {- 1}) = (\mathrm{Ad} \circ g \circ p) \circ (f ^ {- 1}. d f) - (\mathrm{Ad} \circ g \circ p) \circ (g ^ {- 1}. d g).
\]

若 \(f^{-1}.df = g^{-1}.dg\)，则 \((fg^{-1})^{-1}.d(fg^{-1}) = 0\)，即 \(T_x(fg^{-1}) = 0\) 对所有 \(x \in M\) 成立。这证明了第一个断言。第二个断言由《可微与解析流形》R，5.5.3从它推出。

命题 60. 设 G 是李群（当 K 的特征 \textgreater0 时为有限维），M 是 \(\mathbf{C}^r\) 类流形，\(f\) 是从 M 到 G 的 \(\mathbf{C}^r\) 类映射，\(\alpha\) 是 \(f\) 的左微分。则 \(d\alpha + [\alpha]^2 = 0\)。

设 \(\omega\) 是 G 的左典范微分形式。利用第14号命题51的推论1，有

\[
\begin{array}{r l} d \alpha & = d (f ^ {*} (\omega)) = f ^ {*} (d \omega) = f ^ {*} (- [ \omega ] ^ {2}) \\ & = - [ f ^ {*} (\omega) ] ^ {2} = - [ \alpha ] ^ {2}. \end{array}
\]

